% slideSonnet showcase — teaches the tool, in the tool's own format.
%
% Compile:  slidesonnet sty && latexmk -pdf showcase.tex
% Render:   slidesonnet export showcase.pdf -o showcase.mp4 --engine piper
%
% A two-voice dialog: ALEX asks (overlay step 1), the NARRATOR answers
% (overlay step 2). Each step is its own \ssid page with its own narration —
% and, in the sidecar, its own :voice.
\documentclass[aspectratio=169]{beamer}
\usepackage{slidesonnet}
\usepackage{listings}
\usetheme{Madrid}
\usecolortheme{whale}
\setbeamertemplate{navigation symbols}{}
\lstset{basicstyle=\ttfamily\small,breaklines=true,frame=single,
        backgroundcolor=\color{black!4},keepspaces=true}

\title{slideSonnet}
\subtitle{Narration for your slides --- as editable text}
\author{A self-narrated tour}
\date{}

\begin{document}

\begin{frame}[plain]
  \ssid{title}
  \titlepage
\end{frame}

\begin{frame}{What is this?}
  \ssid<1>{intro-q}
  \ssid<2>{intro-a}
  \onslide<1->{\textbf{Alex:} So what does slideSonnet actually do?}
  \onslide<2->{\medskip\textbf{Narrator:} You bring a finished PDF of your slides.
    slideSonnet turns it into a narrated video --- and the narration you're
    hearing right now lives in a plain text file next to that PDF.}
\end{frame}

\begin{frame}{Why a sidecar file?}
  \ssid<1>{why-q}
  \ssid<2>{why-a}
  \onslide<1->{\textbf{Alex:} Why not just put the narration inside the slides?}
  \onslide<2->{\medskip\textbf{Narrator:} Because the PDF is your frozen artifact.
    Keeping narration in a separate, git-diffable file means you can rewrite a
    line and re-render in seconds --- no recompiling, no re-recording.}
\end{frame}

\begin{frame}[fragile]{Step 1 --- mark your slides}
  \ssid{markers}
  Add the macro to your Beamer source and give every slide a stable id:
  \begin{lstlisting}
\usepackage{slidesonnet}
   ...
   \ssid<1>{euler-setup}    % overlay step 1
   \ssid<2>{euler-trick}    % overlay step 2
   ...
  \end{lstlisting}
  The id is stamped invisibly onto each page --- narration binds to it, not to a
  page number.
\end{frame}

\begin{frame}[fragile]{Step 2 --- scaffold the sidecar}
  \ssid{init}
  One command reads the ids straight out of the PDF:
  \begin{lstlisting}
slidesonnet init deck.pdf
  \end{lstlisting}
  You get a blank \texttt{deck.narration} with one block per slide, ready to
  fill in --- by hand, or by an LLM.
\end{frame}

\begin{frame}[fragile]{The sidecar format}
  \ssid{format}
  \begin{lstlisting}
@euler-setup
:voice narrator
We want the sum of one over n squared. [pause 0.8]

@euler-trick
:pace slow
Watch the denominators. [pause 1] This is the trick.
  \end{lstlisting}
  \texttt{@id} headers, optional \texttt{:voice}/\texttt{:pace}, and
  \texttt{[pause N]} --- the one timing primitive.
\end{frame}

\begin{frame}{Step 3 --- edit and preview}
  \ssid<1>{edit-q}
  \ssid<2>{edit-a}
  \onslide<1->{\textbf{Alex:} Do I have to write that text file by hand?}
  \onslide<2->{\medskip\textbf{Narrator:} You can --- or run
    \texttt{slidesonnet edit deck.pdf} for a local editor. Page through the
    deck, type narration beside each slide, and preview the whole thing. The
    preview bakes in your pauses, so it plays exactly like the final video.}
\end{frame}

\begin{frame}[fragile]{Step 4 --- render}
  \ssid{export}
  \begin{lstlisting}
slidesonnet export deck.pdf -o deck.mp4 --engine piper
  \end{lstlisting}
  \begin{itemize}
    \item \texttt{--engine piper} --- free, local. \texttt{elevenlabs} --- cloud, high quality.
    \item \texttt{--silent} or \texttt{--timing estimate} --- a fast rough cut, no TTS.
    \item Subtitles (\texttt{.srt}/\texttt{.vtt}) are written alongside the video.
  \end{itemize}
\end{frame}

\begin{frame}{Only what changed}
  \ssid<1>{cache-q}
  \ssid<2>{cache-a}
  \onslide<1->{\textbf{Alex:} Doesn't re-rendering get slow?}
  \onslide<2->{\medskip\textbf{Narrator:} Audio is cached by content. Edit one
    block and only that line is re-synthesized; everything else is reused. And
    swapping Piper for ElevenLabs never throws away the audio you already paid
    for.}
\end{frame}

\begin{frame}[fragile]{The whole workflow}
  \ssid{recap}
  \begin{lstlisting}
slidesonnet sty                 # drop the LaTeX macro
slidesonnet init   deck.pdf     # scaffold narration
slidesonnet check  deck.pdf     # reconcile ids
slidesonnet edit   deck.pdf     # write & preview
slidesonnet export deck.pdf -o deck.mp4
  \end{lstlisting}
  Every step is scriptable --- an LLM or a Makefile can drive the whole pipeline.
\end{frame}

\begin{frame}[plain]
  \ssid{closing}
  \centering
  {\Large \textbf{slideSonnet}}\\[0.4em]
  Mark your slides. Write narration. Render.\\[1em]
  \texttt{pip install slidesonnet}
\end{frame}

\end{document}
